\scanpage{473}
\mspage{539}

\begin{equation*}
\tag{34}
\left\{
\begin{gathered}
\diagram{\begin{tikzcd}[ampersand replacement=\&]
S\mathcal M_{\rho,\sigma}\arrow[r,"\sim"]\&
\mathcal M_{\rho,\sigma}\mathbin{\underset{\mathbb{G}_m}{\times}}B'_0
\end{tikzcd}}\\
\text{where}\quad\left\{
\begin{gathered}
\diagram{\begin{tikzcd}[ampersand replacement=\&]
\mathcal M_{\rho,\sigma}\arrow[r,"\chi"]\&\mathbb{G}_m\\
B'_0\arrow[r,"{\delta_{B'_0}=\delta_G|B'_0}"]\&\mathbb{G}_m
\end{tikzcd}}
\end{gathered}\right.
\end{gathered}\right.\quad,\quad\text{Whew !}
\end{equation*}

This formula will serve us as the \uline{definition} of
$S\mathcal M_{\rho,\sigma}$ from now on -- it defines at once
a group law on $S\mathcal M_{\rho,\sigma}$, and we find on
$S\mathcal M_{\rho,\sigma}$ \uline{two} extension structures

\begin{equation*}
\tag{35}
\diagram{\begin{tikzcd}[ampersand replacement=\&,column sep=small]
1\arrow[r]\&L_0\arrow[r]\arrow[d,phantom,"\Vert"]\&
S\mathcal M_{\rho,\sigma}\arrow[r]\&\mathcal M_{\rho,\sigma}\arrow[r]\&1\,,\\
\&\operatorname{Ker}(B'_0\xrightarrow{\delta_{B'_0}}\mathbb{G}_m)\&
\widetilde x=(x,\unsure{\underline u|y})\arrow[r,mapsto]\&x
\end{tikzcd}}
\end{equation*}

\begin{equation*}
\tag{36}
\diagram{\begin{tikzcd}[ampersand replacement=\&,column sep=small]
1\arrow[r]\&\mathcal M^+_{\rho,\sigma}\arrow[r]\arrow[d,phantom,"\Vert"]\&
S\mathcal M_{\rho,\sigma}\arrow[r]\&B'_0\arrow[r]\&1\\
\&\operatorname{Ker}(\mathcal M_{\rho,\sigma}\xrightarrow{\chi}\mathbb{G}_m)\&
\begin{gathered}\widetilde x\\\Vert\\(x,\underline u)\end{gathered}
\arrow[r,mapsto]\&\underline u
\end{tikzcd}}
\end{equation*}

The subgroup $S\mathcal M_{\rho,\sigma}[0]$ of
$S\mathcal M_{\rho,\sigma}$ corresponds to the condition $f=1$ i.e.
$\begin{gathered}\underline U\\\Vert\\\Theta_G(x)\end{gathered}=\underline u$
(which already implies
$\begin{gathered}\det\underline U\\\Vert\\\chi(x)\end{gathered}=\det\underline u$
-- it is also the subgroup of the $(x,u)$ in
$\mathcal M_{\rho,\sigma}\times B'_0$ such that we have
$u=\Theta_G(x)$. Thus $S\mathcal M_{\rho,\sigma}[0]$ is the
subgroup of $S\mathcal M_{\rho,\sigma}$ which is the
\inserted{isomorphic} image of $\mathcal M_{\rho,\sigma}[0]$
\deleted{\unsure{under}}
\inserted{(cf. (\unsure{14|19}) page 536')}
under $x\longmapsto(x,\Theta_G(x))$ :

\deleted{\diagram{\begin{tikzcd}[ampersand replacement=\&]
x\arrow[r,mapsto]\&\unsure{(}
\end{tikzcd}}}

\begin{equation*}
\tag{37}
S\mathcal M_{\rho,\sigma}[0]=\operatorname{Im}\left(
\diagram{\begin{tikzcd}[ampersand replacement=\&]
\mathcal M_{\rho,\sigma}[0]\arrow[r,hook]\&S\mathcal M_{\rho,\sigma}\\
x\arrow[r,mapsto]\&(x,\Theta_G(x))
\end{tikzcd}}\right)
\end{equation*}

\[
(\text{where}\quad\mathcal M_{\rho,\sigma}[0]
=\Theta_G^{-1}(B'_0)\quad)
\]

\mspage{472}

\scanpage{474}

Giving a splitting of (35) amounts to a splitting of
\[
\diagram{\begin{tikzcd}[ampersand replacement=\&]
1\arrow[r]\&L_0\arrow[r]\&B'_0\arrow[r]\&\mathbb{G}_m\arrow[r]\&1\,,
\end{tikzcd}}
\]
now there is one ready-made via the subgroup $T_0$. So we deduce from this
an embedding of $\mathcal M_{\rho,\sigma}$ into
$S\mathcal M_{\rho,\sigma}$, which thereby does indeed become a
$1/2$-direct product.

\begin{equation*}
\tag{38}
S\mathcal M_{\rho,\sigma}\simeq\mathcal M_{\rho,\sigma}\cdot L_0
\end{equation*}

$\mathcal M_{\rho,\sigma}$ \unsure{operating} moreover on $L_0$ via
the character $\chi:\mathcal M_{\rho,\sigma}\longrightarrow \mathbb{G}_m$ --
but this description now seems to me much less
``telling'' than (34), for it does not bring out the canonical
homomorphism
\[
\diagram{\begin{tikzcd}[ampersand replacement=\&]
S\mathcal M_{\rho,\sigma}\arrow[r]\&B'_0
\end{tikzcd}}
\]
which is nevertheless the essential thing for the introduction of
$S\mathcal M_{\rho,\sigma}$. Once this is in place, we find
(``for the record'') the

\uline{Theorem}\quad We have a canonical isomorphism

\begin{equation*}
\tag{39}
\diagram{\begin{tikzcd}[ampersand replacement=\&]
S\mathcal M_{\rho,\sigma}\arrow[r]\&
S\deleted{\unsure{\mathcal M}}\inserted{\Sigma}_{\rho,\sigma}^{*}
\end{tikzcd}}
\quad
\begin{gathered}
\text{group}\\
\text{scheme of the}\\
\text{systems }(\alpha,\beta,f)\\
\text{satisfying }(17)\\
\text{and }(31)
\end{gathered}
\end{equation*}

\deleted{\diagram{\begin{tikzcd}[ampersand replacement=\&]
(x,\underline u)\arrow[r,mapsto]\&
(\alpha=\underline u\unsure{\underline a|\Theta_\rho}^{-1},
\beta=\underline u\underline b^{-1},f=\underline u\underline U^{-1})
\end{tikzcd}}}

\deleted{where $\underline a=\Theta_\rho(x)$,
$\underline b=\Theta_\sigma(x)$, $\underline U=\Theta_G(x)$}

\mspage{473}

\scanpage{475}
\mspage{540}

given by :

\begin{equation*}
\tag{40}
\begin{gathered}
\diagram{\begin{tikzcd}[ampersand replacement=\&]
(x,\underline u)\arrow[r,mapsto]\&(\alpha,\beta,f)
\end{tikzcd}}
\quad\text{with}\quad
\alpha=\underline u\underline a^{-1},\quad
\beta=\underline u\underline b^{-1},\quad
f=\underline u\underline U^{-1}\\
\hfill\text{(how beautiful it is !)}\\
\text{where}\quad\underline a=\unsure{\Theta_\rho}(x),\quad
\underline b=\Theta_\sigma(x),\quad\underline U=\Theta_G(x)
\end{gathered}
\end{equation*}

\deleted{\unsure{The} transformation}\inserted{The operation of $L_0$}
on $\Sigma_{\rho,\sigma}$ by

\begin{equation*}
\tag{41}
\diagram{\begin{tikzcd}[ampersand replacement=\&]
(\alpha,\beta,f)\arrow[r]\&\unsure{''|(}\ell\alpha,\ell\beta,\ell f
\end{tikzcd}}
\end{equation*}

corresponds to the operation of \inserted{left} translation by the
subgroup $L_0$ of $S\mathcal M_{\rho,\sigma}$ (cf. (35)), consisting
of the $(1,\ell)$ with $\ell$ in $L_0$, so we recover
$\mathcal M_{\rho,\sigma}\simeq S\mathcal M_{\rho,\sigma}/L_0$
as the quotient \inserted{(group)} scheme of
$S\Sigma_{\rho,\sigma}^{*}$ by the action (41) of $L_0$.

\uline{Remark}\quad In this description of the $(\alpha,\beta,f)$
the new quantities $\underline u$, $\underline a$,
$\underline b$, $\underline U$ occur, which leads to considering the morphism

\begin{equation*}
\tag{42}
\left\{
\diagram{\begin{tikzcd}[ampersand replacement=\&]
S\mathcal M_{\rho,\sigma}\arrow[r]\&B'_0\times T_\rho\times\unsure{N'_\sigma}\times G\\
(x,\underline u)\arrow[r,mapsto]\&(\underline u,\underline a,\underline b,\underline U)
\end{tikzcd}}\right.
\quad\text{with}\quad
\begin{gathered}
\underline a=\Theta_\rho(x)\\
\underline b=\Theta_\sigma(x)\\
\underline U=\Theta_G(x)
\end{gathered}
\end{equation*}

We have in this connection :

\uline{Proposition}\quad The homomorphism (42) is a monomorphism, its image consists of
the systems satisfying

\begin{equation*}
\tag{43}
\det\underline u=\det\underline b=\det\underline U\quad\text{i.e. }\mu,
\quad\text{and}\quad\det\underline a/\mu\in\Gamma\mu_2
\end{equation*}

This is an immediate consequence of the corollary to the proposition on page 536', and of the
definition of $S\mathcal M_{\rho,\sigma}$ as a fiber product.

\mspage{474}

\scanpage{476}

10\textsuperscript{o}) The section-subgroups
$\mathcal M_{\rho,\sigma}(\sigma)$, $\mathcal M_{\rho,\sigma}(\rho)$.

We define $\mathcal M_{\rho,\sigma}(\sigma)$ by the condition

\[
\beta_0=1\quad\text{i.e.}\quad\beta=f\quad\text{i.e. cf. }\unsure{(40)})
\]

\begin{equation*}
\tag{43}
\underline U=\underline b
\end{equation*}

It is evident from the proposition on page 536' that the subgroup

\begin{equation*}
\tag{44}
\diagram{\begin{tikzcd}[ampersand replacement=\&]
\mathcal M_{\rho,\sigma}(\sigma)\arrow[d,phantom,"\Vert\,\unsure{\mathrm{df}}"]
\subset\mathcal M_{\rho,\sigma})\\
\bigl\{(\underline a,\underline b,\underline U)\in\Gamma\mathcal M_{\rho,\sigma}\mid U=\underline b\bigr\}
\end{tikzcd}}
\end{equation*}

provides a section -- when $\underline a,\underline b$ is given
\inserted{in $\mathcal T_{\rho,\sigma}$},
\inserted{i.e.} satisfying the condition

\[
\unsure{\delta_\rho}(\underline a)/\delta_\sigma(\underline b)\in\Gamma\mu_2,
\]

and the simplest way of lifting it to an element
\inserted{$(\underline a,\underline b,\underline U)$}
of $\mathcal M_{\rho,\sigma}$ isn't it by
$(\underline a,\underline b,\underline U\,\unsure{\mathrm{df}}=\underline b)$ !
(And we \unsure{will have} indeed $\delta_G(\underline U)=\delta_\sigma(\underline b)$ .).
We likewise define $\mathcal M_{\rho,\sigma}(\rho)$
by $\alpha_0=1$ i.e. $\alpha=f$ i.e.

\begin{equation*}
\tag{45}
\diagram{\begin{tikzcd}[ampersand replacement=\&]
\mathcal M_{\rho,\sigma}(\rho)\arrow[d,phantom,"\Vert"]\subset\mathcal M_{\rho,\sigma}\\
\bigl\{(\underline a,\underline b,\underline U)\text{ \unsure{in} }\mathcal M_{\rho,\sigma}\mid\underline U=\underline a\bigr\}
\end{tikzcd}}
\end{equation*}

which gives a partial section \inserted{over} the subgroup
$\mathcal T'_{\rho,\sigma}$ of $\mathcal T_{\rho,\sigma}$ defined as

\begin{equation*}
\tag{46}
\diagram{\begin{tikzcd}[ampersand replacement=\&]
\mathcal T'_{\rho,\sigma}\arrow[d,phantom,"\Vert\,\unsure{\mathrm{df}}"]\unsure{\in|\subset}\mathcal T_{\rho,\sigma}\\
\bigl((\underline a,\underline b)\in T_\rho\times N'_\sigma\mid\unsure{\delta_\rho}(\underline a)=\delta_{\unsure{\sigma|\rho}}(\underline b)\bigr)
\arrow[d,phantom,"\Vert"]\\
\operatorname{Ker}\bigl(\mathcal T_{\rho,\sigma}\xrightarrow{\epsilon_3}\mu_2\bigr)
\end{tikzcd}}
\end{equation*}

so

\begin{equation*}
\tag{47}
\mathcal T'_{\rho,\sigma}=T_\rho\times_{\mathbb{G}_m}\unsure{N'_\sigma},\quad\text{and we have}
\end{equation*}

\mspage{475}

\scanpage{477}
\mspage{541}

so

\begin{equation*}
\tag{48}
\diagram{\begin{tikzcd}[ampersand replacement=\&]
\mathcal T'_{\rho,\sigma}/\mathcal T^0_{\rho,\sigma}\arrow[r,"\epsilon_2"]\&\mu_2
\end{tikzcd}}\,.
\end{equation*}

\uline{Remark}\quad These subgroups can be defined, and provide a
section resp. a section \deleted{on}\inserted{over}
$\mathcal T'_{\rho,\sigma}$, as soon as we have (in the general situation
with $\unsure{7^{\circ}|\mathfrak f}$, page 526' ff) embeddings
\diagram{\begin{tikzcd}[ampersand replacement=\&]
T_\rho\arrow[r,hook]\&G
\end{tikzcd}},
\diagram{\begin{tikzcd}[ampersand replacement=\&]
N'_\sigma\arrow[r,hook]\&G
\end{tikzcd}}, such that
$\delta_\rho,\delta_\sigma$ are induced by $\delta_G$.

In the modular case, we also write

\begin{equation*}
\tag{49}
\mathcal M_{\rho,\sigma}(1/2)=\mathcal M_{\rho,\sigma}(\sigma),\quad
\mathcal M_{\rho,\sigma}(-\bar j)=\mathcal M_{\rho,\sigma}(\rho),
\end{equation*}

and we define more \unsure{general\textsuperscript{ly}|generally}

\begin{equation*}
\tag{50}
\mathcal M_{\rho,\sigma}(q),\qquad
q\in\{\unsure{0},1,\infty\};1/2,\unsure{\underline 2|2},-1\}
\end{equation*}

starting from $\mathcal M_{\rho,\sigma}(0)$ and
$\mathcal M_{\rho,\sigma}(1/2)$ by

\begin{equation*}
\tag{51}
\operatorname{int}(\rho_G)\mathcal M_{\rho,\sigma}(q)
=\mathcal M_{\rho,\sigma}(\dot\rho(q))
\end{equation*}

\margin{$\rho_G$ is $\rho$ as an element of the subgroup
$\unsure{\mathcal G}_{\rho,\sigma}^{+}\simeq SG$ of
$\underset{\unsure{\sim}}{\mathcal M}_{\rho,\sigma}$}

$\dot\rho\deleted{(q)}$ is the homographic transformation

\begin{equation*}
\tag{52}
\dot\rho:\quad
\diagram{\begin{tikzcd}[ampersand replacement=\&]
q\arrow[r,mapsto]\&\dfrac{1}{1-q}
\end{tikzcd}}
\end{equation*}

(whose third iterate is the identity, just like the third iterate
of $\widetilde\rho$).

I would like : to define $\mathcal M_{\rho,\sigma}(-j)$, and
\deleted{\unsure{as}} we have a choice within a well-determined
$\underline{\unsure{\mathrm{Sl}|\mathrm{SL}}(2,\mathbb Z)}^{\wedge}$-conjugacy class,
a choice more

\mspage{476.}

\scanpage{478}
\unsure{restricted} would be between
\[
\operatorname{int}(\sigma_i)(\mathcal M_{\rho,\sigma}(0))
\qquad i\in\{0,1,\infty\}
\tag{53}
\]
\mspage{477}